% GENERATED FILE. Edit canonical lessons, then run npm run generate:books.
% canonical-source-hash: fnv1a64:a181df23355f5d24
% canonical-lessons: SA-C266-market-recall, SA-C266-krinami, SA-C266-ask-and-hear-price, SA-C266-complete-transaction, SA-C266-mixed-market-practice, SA-C266-market-transaction-retrieval

\chapter{A small market exchange}
\label{ch:sa-small-market-exchange}
\hlchaptermodality{266}

\begin{chapteropening}
\textbf{By the end of this chapter:} \emph{I can choose a familiar item, ask its price, request it, and close politely.}

Complete chapter 266 by independently choosing a familiar item and completing a small market transaction.
\end{chapteropening}

\section[Practice]{Bring the market words back together}

\label{lesson:SA-C266-market-recall}



\begin{quote}
Say “a stake,” then choose the correct two-person object form, then say the request form “please give.” (\emph{Paṇaḥ; the taught dual object form; dadātu.})
\end{quote}

\subsection*{Your turn}
\begin{itemize}
\raggedright
  \item \emph{Recall:} buys; sells — \emph{krīṇāti; vikrīṇāti}
  \item \emph{Recall:} money; a price; valuable — \textbf{\sk{धनम्}; \sk{मूल्यम्}; \sk{मूल्यवान्}} — \emph{dhanam; mūlyam; mūlyavān}
  \item \emph{Recall:} a shop; a market — \textbf{\sk{आपणः}; \sk{विपणिः}} — \emph{āpaṇaḥ; vipaṇiḥ}
  \item \emph{Recall:} goods for sale; a shopkeeper — \textbf{\sk{पण्यम्}; \sk{आपणिकः}} — \emph{paṇyam; āpaṇikaḥ}
  \item \emph{Recall:} a rupee — \emph{rūpyakam}
\end{itemize}

\begin{grammarlens}[title={What you've built}]
You already own the people, places, actions, and value words. The next step is to make them carry one whole exchange.
\end{grammarlens}

\begin{tcolorbox}[breakable,colback=teal!4,colframe=teal!35!black,title={Before you move on}]
Say four pieces without looking: buys, a price, a shop, a shopkeeper. (\emph{Krīṇāti; mūlyam; āpaṇaḥ; āpaṇikaḥ.})
\end{tcolorbox}
\section[\texorpdfstring{\sk{क्रीणामि} (krīṇāmi)}{क्रीणामि (krīṇāmi)}]{\sk{क्रीणामि} (krīṇāmi) — I buy}

\label{lesson:SA-C266-krinami}



\begin{quote}
Say “a fee,” then “us” as an object, then make one complete polite request for a familiar item. (\emph{Śulkam; asmān; the frame ending in dadātu.})
\end{quote}

\begin{grammarlens}[title={Grammar Lens: the buyer speaks}]
You know \textbf{\sk{क्रीणाति}} (\emph{krīṇāti}), “he or she buys.” When the buyer speaks for themselves, keep the whole new form together:

\begin{quote}
\textbf{\sk{अहं} \sk{फलं} \sk{क्रीणामि।}}
\end{quote}

\begin{quote}
\emph{Ahaṃ phalaṃ krīṇāmi.}
\end{quote}

\begin{quote}
I buy fruit.
\end{quote}

Here \textbf{\sk{क्रीणामि}} is the \textbf{first-person singular present}. The final \textbf{-\sk{मि}} (\emph{-mi}) points to the speaker, as it does in forms you already know. Both checked-in A1 mocks use this complete form. This is \textbf{not a complete present paradigm}: the other persons, numbers, and class-nine forms wait for later.
\end{grammarlens}

\subsection*{Your turn}
\begin{itemize}
\raggedright
  \item \emph{Contrast:} buys; I buy — \emph{krīṇāti; krīṇāmi}
  \item \emph{Say it:} “I buy fruit” — \emph{ahaṃ phalaṃ krīṇāmi}
  \item \emph{Say it:} the buyer form once more — \emph{krīṇāmi}
\end{itemize}

\begin{tcolorbox}[breakable,colback=teal!4,colframe=teal!35!black,title={Before you move on}]
Change \textbf{\sk{क्रीणाति}}, “buys,” into the one taught speaker form. (\textbf{\sk{क्रीणामि}}.)
\end{tcolorbox}
\section[(ask and hear one price)]{Ask the price and hear the answer}

\label{lesson:SA-C266-ask-and-hear-price}



\begin{quote}
Say “news,” then “you two” as an object. (\emph{Vārtā; yuṣmān.})
\end{quote}

\begin{grammarlens}[title={Grammar Lens: one price turn}]
Put two known words together to ask the question used in the A1 mocks:

\begin{quote}
\textbf{\sk{मूल्यं} \sk{कियत्}?}
\end{quote}

\begin{quote}
\emph{Mūlyaṃ kiyat?}
\end{quote}

\begin{quote}
How much is the price?
\end{quote}

One short answer is:

\begin{quote}
\textbf{\sk{विंशतिः} \sk{रूप्यकाणि।}}
\end{quote}

\begin{quote}
\emph{Viṃśatiḥ rūpyakāṇi.}
\end{quote}

\begin{quote}
Twenty rupees.
\end{quote}

You already know \textbf{\sk{रूप्यकम्}}, one rupee. Keep \textbf{\sk{रूप्यकाणि}} together inside this useful price answer. This small exchange does not teach the complete neuter plural system or every way to state a price.
\end{grammarlens}

\subsection*{Your turn}
Say these aloud:

\begin{itemize}
\raggedright
  \item ask the price — \emph{mūlyaṃ kiyat?}
  \item answer “twenty rupees” — \emph{viṃśatiḥ rūpyakāṇi}
  \item the question and answer as one turn
\end{itemize}

\begin{tcolorbox}[breakable,colback=teal!4,colframe=teal!35!black,title={Before you move on}]
Ask one price, then answer with twenty rupees.
\end{tcolorbox}
\section[Practice]{Complete one small transaction}

\label{lesson:SA-C266-complete-transaction}



\begin{quote}
Say “a drama,” then explain how the taught object forms change for person and number. (\emph{Nāṭakam; use the object-person-and-number contrast.})
\end{quote}

\subsection*{Your turn: five turns}
Build the exchange one complete turn at a time:

\begin{enumerate}
\raggedright
  \item \textbf{\sk{अहं} \sk{फलं} \sk{क्रीणामि।}} — \emph{Ahaṃ phalaṃ krīṇāmi.} — I buy fruit.
  \item \textbf{\sk{मूल्यं} \sk{कियत्}?} — \emph{Mūlyaṃ kiyat?} — How much is the price?
  \item \textbf{\sk{विंशतिः} \sk{रूप्यकाणि।}} — \emph{Viṃśatiḥ rūpyakāṇi.} — Twenty rupees.
  \item \textbf{\sk{कृपया} \sk{मह्यं} \sk{फलं} \sk{ददातु।}} — \emph{Kṛpayā mahyaṃ phalaṃ dadātu.} — Please give me fruit.
  \item \textbf{\sk{धन्यवादः।}} — \emph{Dhanyavādaḥ.} — Thank you.
\end{enumerate}

Say the five turns once without breaking the exchange apart.

\begin{grammarlens}[title={What you've built}]
The vocabulary was already yours. Now it carries one complete A1 transaction.
\end{grammarlens}

\begin{tcolorbox}[breakable,colback=teal!4,colframe=teal!35!black,title={Before you move on}]
Name the five jobs in order: choose and buy, ask the price, hear the price, request the item, close politely.
\end{tcolorbox}
\section[Practice]{Choose and complete the exchange}

\label{lesson:SA-C266-mixed-market-practice}



\begin{quote}
Say “a poem.” (\emph{Kāvyam.})
\end{quote}

\subsection*{Your turn: independent round one}
No Sanskrit model is shown before you answer.

You are speaking to a shopkeeper at a market. Choose one familiar item. Say that you buy it. Ask the price. The shopkeeper says that the price is twenty rupees. Request the item politely, then close with thanks.

\begin{grammarlens}[title={What you've built: check after speaking}]
One complete answer with fruit is \textbf{\sk{अहं} \sk{फलं} \sk{क्रीणामि।} \sk{मूल्यं} \sk{कियत्}? \sk{विंशतिः} \sk{रूप्यकाणि।} \sk{कृपया} \sk{मह्यं} \sk{फलं} \sk{ददातु।} \sk{धन्यवादः।}} Your item may differ; the five jobs must remain.
\end{grammarlens}

\begin{tcolorbox}[breakable,colback=teal!4,colframe=teal!35!black,title={Before you move on}]
Which two new pieces held the old vocabulary together? (The buyer form \textbf{\sk{क्रीणामि}} and the price-exchange pattern.)
\end{tcolorbox}
\section[Practice]{Retrieve the whole market transaction}

\label{lesson:SA-C266-market-transaction-retrieval}



\begin{quote}
Say “a film,” then “I will go.” (\emph{Citrapaṭaḥ; gamiṣyāmi.})
\end{quote}

\subsection*{Your turn: independent round two}
No Sanskrit model is shown before you answer.

Choose a shop or market and a different familiar item. Say that you buy it. Ask its price. The shopkeeper says that the price is twenty rupees. Request the item politely and close with thanks.

\begin{grammarlens}[title={What you've built: check after speaking}]
A complete answer still has five jobs: say what you buy; ask the price; hear the rupee amount; request the item; thank the shopkeeper. Your chosen item may differ from the first round.
\end{grammarlens}

\begin{tcolorbox}[breakable,colback=teal!4,colframe=teal!35!black,title={Before you move on}]
You can now complete one small A1 market exchange with familiar language. This is \textbf{not} complete commerce, currency, bargaining, consumer-rights, or present paradigm coverage, and it is not evidence of exam readiness or learner success.
\end{tcolorbox}
